\section*{ABSTRACT}
\phantomsection\addcontentsline{toc}{section}{\numberline{}ABSTRACT}

This report presents a comparative study of Orthogonal Frequency Division
Multiplexing (OFDM) and Single-Carrier Frequency Division Multiple Access
(SC-FDMA), two important multicarrier transmission techniques widely used
in modern wireless communication systems. OFDM is well known for its high
spectral efficiency, robustness against multipath fading, and efficient
frequency-domain equalization. However, one major limitation of OFDM is
its high Peak-to-Average Power Ratio (PAPR), which reduces power amplifier
efficiency and introduces nonlinear distortion in practical transmission
systems.

SC-FDMA was introduced to address the high-PAPR problem of OFDM by
applying a Discrete Fourier Transform (DFT) spreading stage prior to OFDM
modulation. This additional precoding operation allows SC-FDMA to retain
many advantages of OFDM while exhibiting single-carrier-like transmission
characteristics and lower PAPR. Consequently, SC-FDMA is more suitable for
uplink communication in battery-powered mobile devices, where transmitter
power efficiency is critical.

In this study, the operating principles of OFDM and SC-FDMA are analyzed
theoretically and evaluated through MATLAB-based simulations. The systems
are implemented using 16-QAM modulation over multipath Rayleigh fading
channels with Zero Forcing (ZF) equalization. Performance evaluation is
conducted using Bit Error Rate (BER) and Complementary Cumulative
Distribution Function (CCDF) analysis of PAPR. Simulation results show
that both systems achieve comparable BER performance under identical
channel conditions, while SC-FDMA demonstrates significantly lower PAPR
than OFDM.

The results highlight the tradeoff between transmission performance and
power efficiency, and explain the practical adoption of OFDM in downlink
systems and SC-FDMA in LTE uplink communications.

\newpage
\clearpage